\chapter{Thắng thua trong lòng mình}
\markboth{Thắng thua trong lòng mình}{Inner Victory and Loss}

\section{Một người kiếm được rất nhiều nhưng luôn thấy thiếu.}
\begin{truyen}{Một người kiếm được rất nhiều nhưng luôn thấy thiếu.}{If enough never comes, you are still losing.}
\chuhoa{A}{nh đạt những cột mốc người khác mơ ước, nhưng cứ xong một đích lại hoảng vì phải có thêm. Nhìn từ ngoài, anh đang thắng; nhìn từ trong, lòng anh chưa từng thôi thua trước cảm giác không đủ.}
\textit{(He reached milestones others dreamed of, but after each one came more anxiety for the next. From outside, he was winning; inside, he kept losing to a sense of never being enough.)}
\end{truyen}
\begin{baihoc}
Chiến thắng không làm dịu được lòng tham thì chưa phải là chiến thắng trọn vẹn.
\textit{(If victory cannot quiet greed, it is not a complete victory.)}
\end{baihoc}
\begin{ghinhoanh}\textbf{Short English to remember:}

If enough never comes, you are still losing.
\end{ghinhoanh}
\begin{tuhocnhanh}
1. Em có đang chạy mà không biết điểm đủ của mình không? \textit{(Are you running without knowing your point of enough?)}

2. Điều gì làm em thấy đã đủ? \textit{(What truly makes you feel enough?)}
\end{tuhocnhanh}
\ngancach

\section{Một người chịu nhường nhưng ngủ rất yên.}
\begin{truyen}{Một người chịu nhường nhưng ngủ rất yên.}{Peace can be a better win.}
\chuhoa{C}{hị biết mình có thể cãi tới cùng, nhưng chọn dừng vì thấy câu chuyện không đáng. Người ngoài bảo chị thua, còn chị thấy lòng nhẹ. Hóa ra có những cái thua ở miệng lại là thắng ở giấc ngủ.}
\textit{(She knew she could argue to the end, but chose to stop because the matter was not worth it. Others said she lost; she felt light inside. Some losses in speech become victories in sleep.)}
\end{truyen}
\begin{baihoc}
Bình yên đôi khi là phần thưởng của việc không cố thắng cho bằng được.
\textit{(Peace is sometimes the reward for not forcing a victory.)}
\end{baihoc}
\begin{ghinhoanh}\textbf{Short English to remember:}

Peace can be a better win.
\end{ghinhoanh}
\begin{tuhocnhanh}
1. Có việc gì em nên dừng tranh hơn là tiếp tục hơn thua? \textit{(What should you stop arguing over instead of continuing to compete?)}

2. Em đã bao giờ thấy nhẹ vì chịu nhường chưa? \textit{(Have you ever felt relieved by letting go of a win?)}
\end{tuhocnhanh}
\ngancach

\section{Một người không được vinh danh nhưng vẫn làm việc tử tế.}
\begin{truyen}{Một người không được vinh danh nhưng vẫn làm việc tử tế.}{Character does not need a podium.}
\chuhoa{A}{nh không được nêu tên trong buổi tổng kết, dù đóng góp không ít. Anh buồn nhưng không cay cú, vẫn làm việc chỉnh chu như cũ. Chính thái độ đó làm người khác nể anh hơn những tấm bảng khen.}
\textit{(He was not named in the final recognition event despite contributing much. He felt sad but not bitter, and continued to work carefully. That attitude earned more respect than certificates could.)}
\end{truyen}
\begin{baihoc}
Thắng lớn nhất đôi khi là không để sự bỏ sót của người khác làm mình biến chất.
\textit{(The greatest victory is sometimes not letting other people's omission corrupt your character.)}
\end{baihoc}
\begin{ghinhoanh}\textbf{Short English to remember:}

Character does not need a podium.
\end{ghinhoanh}
\begin{tuhocnhanh}
1. Khi không được ghi nhận, em thay đổi thái độ ra sao? \textit{(How do you change when you are not recognized?)}

2. Điều gì giữ em làm đúng dù không ai vỗ tay? \textit{(What keeps you doing right even without applause?)}
\end{tuhocnhanh}
\ngancach

\section{Một người thắng người khác nhưng không thắng được cơn nóng giận.}
\begin{truyen}{Một người thắng người khác nhưng không thắng được cơn nóng giận.}{Outer victory cannot hide inner defeat.}
\chuhoa{S}{au cuộc cãi vã, anh biết mình lý lẽ hơn. Nhưng suốt đêm anh vẫn bực, vẫn dựng lại từng câu trong đầu. Anh nhận ra mình thắng người đối diện mà thua chính cảm xúc của mình.}
\textit{(After the argument, he knew his reasoning was stronger. Yet all night he replayed it in anger. He realized he had defeated the other person while losing to his own emotions.)}
\end{truyen}
\begin{baihoc}
Không làm chủ được lòng mình thì cái thắng bên ngoài vẫn còn dang dở.
\textit{(If you cannot govern your inner self, outer victory remains unfinished.)}
\end{baihoc}
\begin{ghinhoanh}\textbf{Short English to remember:}

Outer victory can hide inner defeat.
\end{ghinhoanh}
\begin{tuhocnhanh}
1. Điều gì thường thắng được em trong lòng? \textit{(What usually defeats you inside?)}

2. Em cần học điều gì để bớt bị cảm xúc kéo đi? \textit{(What do you need to learn so emotions pull you less?)}
\end{tuhocnhanh}
\ngancach

\section{Một người chấp nhận mình chưa bằng ai đó và bớt ghen.}
\begin{truyen}{Một người chấp nhận mình chưa bằng ai đó và bớt ghen.}{Honest acceptance can free the heart.}
\chuhoa{K}{hi thôi giả vờ mình không quan tâm, cô bắt đầu học từ người giỏi hơn thay vì ghen với họ. Từ đó, lòng cô nhẹ hơn và tiến bộ nhanh hơn. Có những lần nhận mình chưa giỏi lại là lúc mình bắt đầu thắng thật.}
\textit{(When she stopped pretending not to care, she began learning from those ahead instead of envying them. Her heart became lighter and her progress faster. Admitting she was not yet good enough was the start of real victory.)}
\end{truyen}
\begin{baihoc}
Thừa nhận mình còn kém không làm ta nhỏ đi; nó mở đường để ta lớn lên.
\textit{(Admitting you are still behind does not make you smaller; it opens a path to growth.)}
\end{baihoc}
\begin{ghinhoanh}\textbf{Short English to remember:}

Honest acceptance can free the heart.
\end{ghinhoanh}
\begin{tuhocnhanh}
1. Em đang ghen hay đang học từ người giỏi hơn? \textit{(Are you envying or learning from those ahead of you?)}

2. Điều gì em cần thừa nhận thật với bản thân? \textit{(What do you need to admit honestly to yourself?)}
\end{tuhocnhanh}
\ngancach